\section{A uniform bound for pencils}\label{sec:pencil}

The preceding construction tests a finite collection of functions at a
finite Kummer level.  We now combine such tests with five fixed members
of a pencil.  Riemann--Hurwitz bounds the degree of every residual
pencil, which in turn bounds the support of the entire original pencil.
Bounded-genus curve sections and Riemann--Hurwitz already underlie the
finite-support argument of Bogomolov and Tschinkel
\cite[Section~6, proof of Proposition~6.1]{BT2011}.
Here the fixed anchors give one bound for every member of the pencil,
and the finite residual tests allow arbitrary algebraically closed constants.

Keep the model $X$ of the preceding section and the Kummer isomorphism
$\Theta:\widehat K\to\widehat L$.
For $h\in\widehat L$ and $e\ge1$, define
\[
 S_e(h)=\{D\subset X\text{ prime divisor}:
                \operatorname{ord}_D(h)\not\equiv0\pmod{\ell^e}\},
 \qquad
 \operatorname{Supp}_X(h)=\{D:\operatorname{ord}_D(h)\ne0\}.
\]
The order maps extend continuously to $\widehat L$.
Each $S_e(h)$ is finite, because $h$ modulo $\ell^e$ is represented by an
ordinary rational function.  The exact support need not be finite a
priori.  The sets $S_e(h)$ increase with $e$, and their union is
$\operatorname{Supp}_X(h)$.

\begin{theorem}[Uniform pencil support]\label{pencil:uniform}
For every $t\in K\setminus k$, there is a finite constant $B_t$ such that
\[
 \sum_{D\in\operatorname{Supp}_X(\Theta[t-a])}\deg D\le B_t
 \qquad\text{for every }a\in k.
\]
In particular, all these supports are finite, with a bound independent
of the member of the pencil.
\end{theorem}

\begin{proof}
Choose five distinct anchors $a_1,\ldots,a_5\in k$, including $0$.
In characteristic zero include $1$ and $\ell$ as well.
Put $h_a=\Theta[t-a]$ and
\[
 N_i=\sum_{D\in S_1(h_{a_i})}\deg D\quad(1\le i\le5),
 \qquad c_\ell=5(1-1/\ell)-2=3-5/\ell.
\]
The $N_i$ are finite before any exact-support assertion is known.
Moreover $c_\ell\ge1/2$, so the constant
\begin{equation}\label{pencil:constant}
 B=\max\left\{1,\frac{2G_X-2+\sum_{i=1}^5N_i}{c_\ell}\right\}
\end{equation}
is finite and depends only on $X$, $\Theta$, $t$ and the fixed anchors.
We prove the theorem with $B_t=2B$.

Fix one additional $a\in k$, and list the distinct elements among the
five anchors and $a$ as $a_j$.  Write $u_j=t-a_j$.
For all sufficiently large $e$, the classes $[u_j]$ modulo $q=\ell^e$
are nonzero and pairwise distinct and generate a noncyclic subgroup.
Indeed the functions $t-a_j$ are independent modulo constants: a
multiplicative relation would give a constant rational function of the
transcendental element $t$ and hence have all exponents zero.
Their images in the free abelian group $K^\times/k^\times$ are
therefore independent.  Integral evaluations detect nonzero elements
and nonzero two-by-two determinants, which remain nonzero modulo a
sufficiently large power of $\ell$.
Equivalently, one can prolong two distinct point valuations of $k(t)$
to $K$ to obtain a diagonal evaluation matrix with nonzero integer
entries.  The necessary lower bound on $e$ may depend on $a$.

Choose representatives $g_j\in L^\times$ of $h_{a_j}$ modulo $q$,
and apply Lemma~\ref{tests:residual}.
It gives $w$, a residual curve $C$, and the same noncyclic subgroup with
distinct nonzero generators in $(Lw)^\times/(Lw)^{\times q}$.
Evaluation on $D_w/I_w=W_{Lw}$ detects these classes.
By Proposition~\ref{local:transport}, there is a corresponding quasi-prime
$(d-1)$-divisor $v$ of $K/k$, and the evaluations of the $u_j$ on $D_v$
modulo $q$ are likewise nonzero, distinct and noncyclic.

\emph{Normalize the residual pencil.}
We first show that
\begin{equation}\label{pencil:no-relative-value}
 v(k(t))=vk.
\end{equation}
Otherwise $v(k(t))/vk$ is a nonzero subgroup of the finitely generated
free group $vK/vk$.  The valuation transcendence-degree inequality for
$k(t)/k$ makes it cyclic and gives $k(t)v=kv$.
Every function of value in $vk$ then becomes a principal unit after
multiplication by constants.  Since $D_v$ kills constants and principal
units, evaluation on $D_v$ factors through this cyclic relative value
group.  Its image modulo $q$ is cyclic, a contradiction.

All the $v(u_j)$ have a common value $\gamma\in vk$.
For if $v(u_j)<v(u_h)$, then
$u_j/(a_h-a_j)=1+u_h/(a_h-a_j)$ is a principal unit, forcing the
evaluation of $u_j$ to vanish.
Also $v(a_j-a_h)=\gamma$ whenever $j\ne h$:
a strictly larger value would make $u_j/u_h$ a principal unit and
force two evaluations to coincide.
Choose $b\in k^\times$ with $v(b)=\gamma$, and fix one listed anchor
$a_0$.  Set
\[
 \tau_0=\overline{(t-a_0)/b},\qquad
 \beta_j=\overline{(a_j-a_0)/b}.
\]
These expressions are defined, the $\beta_j$ are distinct elements of
$kv$, and
\begin{equation}\label{pencil:residual-linear}
 \overline{u_j/b}=\tau_0-\beta_j.
\end{equation}
The element $\tau_0$ is nonconstant: a constant residue would make
$(t-a_0)/b$ a constant times a principal unit and give zero evaluation.

The residue characteristic is different from $\ell$.
This is automatic if $k$ has positive characteristic.
In characteristic zero, the anchor differences $1-0$ and $\ell-0$
have the same value, so $v(\ell)=v(1)=0$.
Thus $\ell$ is nonzero in $kv$ in that case as well.

\emph{Compare the two residual curves.}
The constants $lw$ are algebraic over a finite field.
Proposition~\ref{local:curves} therefore gives the same assertion for $kv$
and identifies the full families of point inertia lines in
$W_{Lw}$ and $W_{Kv}$.
Let $C'/kv$ be the smooth projective curve with function field $Kv$.
By Lemma~\ref{prelim:curve-kernel}, the mod-$\ell$ kernels of the point
order maps have dimensions $2g(C)$ and $2g(C')$.
The residual isomorphism identifies these kernels, and hence
\begin{equation}\label{pencil:genus}
 g(C')=g(C)\le G_X.
\end{equation}
Generators of matched point inertia lines differ by units of $\Zell$.
Consequently the number of nonzero point orders modulo $\ell$ or $q$
is preserved by the comparison.
Because constants have zero Kummer class, evaluation identifies the
classes of the functions in \eqref{pencil:residual-linear} with those
of $\overline g_j$ modulo $q$.

\emph{Use the five anchors to bound the degree.}
Write $p=\operatorname{char}(kv)>0$, using the finite-constant
recognition above.
Remove inseparability by writing $\tau_0=\tau^{p^r}$ with $r\ge0$
maximal.  Such an $r$ is finite, as a nonzero integer point order of
$\tau_0$ cannot be divisible by arbitrarily high powers of $p$.
The function $\tau$ gives a separable morphism
$\tau:C'\to\mathbb P^1$ of some degree $n\ge1$.
Let $\beta'_j$ be the $p^r$-th root of $\beta_j$ in $kv$.
These constants are distinct, and
\[
 \tau_0-\beta_j=(\tau-\beta'_j)^{p^r}.
\]
Since $p\ne\ell$, this power does not affect whether any point order
is zero modulo $\ell$ or $q$.

Over the value $\beta'_i$ corresponding to anchor $i$, at most $N_i$
points have ramification index not divisible by $\ell$.
Indeed such a point contributes a nonzero order modulo $\ell$ to
$\tau_0-\beta_i$; transport to $C$ and
\eqref{tests:support} give the bound $N_i$.
Every other point of that fiber has ramification index at least
$\ell$, so there are at most $n/\ell$ of them.
If $r_i$ denotes the total number of points in the fiber, then
\[
 r_i\le N_i+n/\ell,\qquad
 \sum_{P\mid\beta'_i}(e_P-1)=n-r_i
       \ge(1-1/\ell)n-N_i.
\]
For a separable morphism of smooth curves, the different exponent at
$P$ is at least $e_P-1$, including in the presence of wild ramification.
The five fibers are disjoint.  Riemann--Hurwitz therefore yields
\[
 2g(C')-2+2n
 \ge\sum_{i=1}^5\sum_{P\mid\beta'_i}(e_P-1)
 \ge5(1-1/\ell)n-\sum_{i=1}^5N_i.
\]
Together with \eqref{pencil:genus} this gives $c_\ell n\le
2G_X-2+\sum_iN_i$, and hence $n\le B$.
This argument uses the different inequality, not the tame
Riemann--Hurwitz formula; see \cite[Section~53.12, Tag~0C1B]{Stacks}.

\emph{Pass from finite tests to exact supports.}
For every tested $j$, all zeros and poles of $\tau-\beta'_j$ lie in two
fibers, over $\beta'_j$ and infinity.  There are at most $2n$ such
points.  The number of its nonzero point orders modulo $q$ is therefore
at most $2B$.
Transport this bound to $C$ and apply \eqref{tests:support} to obtain
\begin{equation}\label{pencil:finite-level-bound}
 \sum_{D\in S_e(h_a)}\deg D\le2B.
\end{equation}
The right side is independent of $a$, $e$, the finite descent data and
the residual test.  For each fixed $a$, the inequality holds for every
sufficiently large $e$.

Finally $S_e(h_a)$ increases to $\operatorname{Supp}_X(h_a)$.
If the exact support had total degree greater than $2B$, a finite
subset would already have degree greater than $2B$; that subset would
lie in $S_e(h_a)$ for all sufficiently large $e$, contradicting
\eqref{pencil:finite-level-bound}.
Every prime divisor has positive integer degree, so the exact support
is finite.  This proves the asserted uniform bound.
\end{proof}

\begin{corollary}\label{pencil:finite}
For every $z\in K^\times$, the class $\Theta[z]$ has finite support on
$X$.
\end{corollary}

\begin{proof}
If $z\notin k$, apply Theorem~\ref{pencil:uniform} to $t=z$ and $a=0$.
If $z\in k^\times$, its Kummer class is zero because $k^\times$ is
divisible.
\end{proof}
